\documentclass[14pt]{extarticle}
\usepackage{natbib}

\usepackage{xcolor}
\usepackage{amsmath}
\newcommand{\tuple}[1]{ \langle #1 \rangle }
%\usepackage{automata}
\usepackage{times}
\usepackage{ltablex}
\usepackage{tasks}
\usepackage{threeparttable}
\usepackage{booktabs}
\usepackage{xcolor} % For custom colors
\usepackage{tikz} % For styling enumerate numbers
\usepackage{tcolorbox} % For colored box styling
\usepackage{amsmath, amsfonts, amssymb, amsthm} % Math related
\usepackage{natbib}
\usepackage{fontspec}
\usepackage{luatexja}
\usepackage[mathscr]{euscript}

%%%%%% Template
\usepackage{hyperref}
\hypersetup{colorlinks=true,allcolors=blue}

\usepackage{vmargin}
\setpapersize{USletter}
\setmarginsrb{1.0in}{1.0in}{1.0in}{0.6in}{0pt}{0pt}{0pt}{0.4in}

% HOW TO USE THE ABOVE:
%\setmarginsrb{leftmargin}{topmargin}{rightmargin}{bottommargin}{headheight}{headsep}{footheight}{footskip}
%\raggedbottom
% paragraphs indent & skip:
\parindent  0.3cm
\parskip    -0.01cm

\usepackage{tikz}
\usetikzlibrary{backgrounds}

% hyphenation:
% \hyphenpenalty=10000 % no hyphen
% \exhyphenpenalty=10000 % no hyphen
\sloppy

% notes-style paragraph spacing and indentation:
\usepackage{parskip}
\setlength{\parindent}{0cm}

% let derivations break across pages
\allowdisplaybreaks

\newcommand{\orange}[1]{\textcolor{orange}{#1}}
\newcommand{\blue}[1]{\textcolor{blue}{#1}}
\newcommand{\red}[1]{\textcolor{red}{#1}}
\newcommand{\freq}[1]{{\bf \sf F}(#1)}
\newcommand{\datafreq}[2]{{{\bf \sf F}_{#1}(#2)}}

\def\qqquad{\quad\qquad}
\def\qqqquad{\qquad\qquad}

%%%%%%%%%%%%%%%%%%%%%%%%%%%%%%%%%%%%%%%%%%%%%%%%%%%%%%%%%%%%%%%%%%%%%%%%%%%%%%%%
%%%%%%%%%%%%%%%%%%%%%%%%%%%%%%%%%%%%%%%%%%%%%%%%%%%%%%%%%%%%%%%%%%%%%%%%%%%%%%%%

% fill-in-blank question style, found in https://tex.stackexchange.com/a/505089

\usepackage{ifthen}
\usepackage{tocloft}
\usepackage{exercise}
% \usepackage{xcolor}

% Set the Show Answers Boolean
\newboolean{showAns}
\setboolean{showAns}{false}
\newcommand{\showAns}{\setboolean{showAns}{true}}

% The length of the Answer line
\newlength{\answerlength}
\newcommand{\anslen}[1]{\settowidth{\answerlength}{#1}}

% ans command that indicates space for an answer or shows the answer in red
\newcommand{\ans}[1]{\settowidth{\answerlength}{\hspace{2ex}#1\hspace{2ex}}%
    \ifthenelse{\boolean{showAns}}%
        {\textcolor{red}{\underline{\hspace{2ex}#1\hspace{2ex}}}}%
        {\underline{\hspace{\answerlength}}}}%

\newcommand{\details}[1]{\settowidth{\answerlength}{#1}%
    \ifthenelse{\boolean{showAns}}%
        {\begin{tcolorbox}[colback=blue!5!white,colframe=blue!75!black,title=Solution] #1 \end{tcolorbox}}%
        {}}%

% Formatting how multiple choices Questions are formated.
\settasks{label=(\Alph*), label-width=30pt}


% Some commands for the Exercise Question package
\renewcommand{\QuestionNB}{\Large\protect\textcircled{\small\bfseries\arabic{Question}}\ }
\renewcommand{\ExerciseHeader}{} %no header
\renewcommand{\QuestionBefore}{3ex} %Space above each Q
\setlength{\QuestionIndent}{8pt} % Indent after Q number


% To create the list of answers with tocloft...
\newcommand{\listanswername}{Answers}
\newlistof[Question]{answer}{Answers}{\listanswername}

% Creates a TOC for Answers
\newcounter{prevQ}
\newcommand{\answer}[1]{\refstepcounter{answer}%
\ans{#1}%
\ifnum\theQuestion=\theprevQ%
        \addcontentsline{Answers}{answer}{\protect\numberline{}#1}% don't include the Q number
        \else%
        \addcontentsline{Answers}{answer}{\protect\numberline{\theQuestion}#1}%
        \setcounter{prevQ}{\value{Question}}%
        \fi%
        }%

% \hyphenpenalty=10000 % no hyphen
% \exhyphenpenalty=10000 % no hyphen
\sloppy              % hyphen

\newcommand{\HRule}{\rule{\linewidth}{0.5mm}}
\newcommand{\Hrule}{\rule{\linewidth}{0.3mm}}

%tocloft formatting listofanswers
\renewcommand{\cftAnswerstitlefont}{\bfseries\large}
\renewcommand{\cftanswerdotsep}{\cftnodots}
\cftpagenumbersoff{answer}
\addtolength{\cftanswernumwidth}{10pt}

\makeatletter% since there's an at-sign (@) in the command name
\renewcommand{\@maketitle}{%
  \parindent=0pt% don't indent paragraphs in the title block
  \centering
  {\Large \bfseries\textsc{\@title}} \\
  \vspace{5pt}
  {\large \textit{\@author}} \\
  \HRule \\
  \vspace{1em}
}
\makeatother% resets the meaning of the at-sign (@)


% --- Fonts (robust fallback) ---
\IfFontExistsTF{Crimson Pro Light}{
  \setmainfont{Crimson Pro Light}[
    ItalicFont={* Italic},
    BoldFont={Crimson Pro Medium},
    BoldItalicFont={Crimson Pro Medium Italic}]
  \setsansfont{Crimson Pro Light}[
    ItalicFont={* Italic},
    BoldFont={Crimson Pro Medium},
    BoldItalicFont={Crimson Pro Medium Italic}]
}{
  \setmainfont{HaranoAjiMincho}
  \setsansfont{HaranoAjiGothic}
}
\title{Problem Set 2}
\author{Macroeconomics II (Spring 2026) \\ Hui-Jun Chen}



% ------------------------------------------------------------
% Explanations for each option (only printed when \showAns)
% Usage: \ExplainChoices{A}{B}{C}{D}
% ------------------------------------------------------------
\newcommand{\ExplainChoices}[4]{%
\ifthenelse{\boolean{showAns}}{%
\vspace{0.3em}\noindent\textcolor{red}{\textbf{Explanations.}
\begin{itemize}
\item[(A)] #1
\item[(B)] #2
\item[(C)] #3
\item[(D)] #4
\end{itemize}}
}{}%
}

\begin{document}

\maketitle

% Uncomment to display answers in red:
% \showAns

\begin{Exercise}

\section*{Instructions}
\begin{itemize}
\item This problem set is based on Karl Whelan's \emph{Advanced Macroeconomics} slides:
\textbf{Slides 1--3} on the IS--MP--PC model and the Taylor principle.
\item 40 multiple-choice questions. Exactly one option is correct for each question.
\end{itemize}

\vspace{0.6em}
\Hrule
\vspace{0.8em}

\section*{Problem 1 (Questions 1--20): Phillips curve, IS, MP, and the IS--MP curve}

\Question \textbf{Original Phillips curve (historical claim).} Phillips (1958) documented a negative relationship between:
\answer{B}
\begin{tasks}(1)
\task price inflation and real GDP growth
\task wage inflation and unemployment
\task unemployment and real interest rates
\task money growth and unemployment
\end{tasks}
\ExplainChoices{Wrong: Phillips’ original evidence was about wage inflation, not price inflation or GDP growth.}{Correct: the original Phillips curve related wage inflation to unemployment (negative relationship).}{Wrong: the relationship was not between unemployment and real interest rates.}{Wrong: money growth is a different relationship (quantity theory), not Phillips (1958).}



\Question \textbf{Friedman's critique.} Friedman (1968) argued that the Phillips curve tradeoff disappears in the long run mainly because:
\answer{C}
\begin{tasks}(1)
\task technology shocks vanish in the long run
\task central banks cannot influence nominal interest rates
\task inflation expectations adjust, shifting the Phillips curve
\task unemployment is always zero in equilibrium
\end{tasks}
\ExplainChoices{Wrong: long-run breakdown is not about shocks vanishing.}{Wrong: central banks can set nominal rates; that’s not the critique.}{Correct: once expectations adjust, the curve shifts; no permanent unemployment–inflation tradeoff.}{Wrong: equilibrium unemployment is not zero in general.}



\Question \textbf{Natural rate language.} The ``natural'' level of output $y_t^\ast$ is defined as the level consistent with:
\answer{D}
\begin{tasks}(1)
\task zero inflation
\task the government budget being balanced
\task zero real interest rates
\task unemployment equal to its natural rate
\end{tasks}
\ExplainChoices{Wrong: zero inflation is not the definition of natural output.}{Wrong: balanced budgets are fiscal, not defining natural output.}{Wrong: natural output does not require zero real rates.}{Correct: natural output corresponds to unemployment at its natural rate (no cyclical unemployment).}



\Question \textbf{Whelan Phillips curve (form).} The Phillips curve used is:
\answer{A}
\begin{tasks}(1)
\task $\pi_t=\pi_t^e+\gamma(y_t-y_t^\ast)+\varepsilon_t^\pi$
\task $\pi_t=\gamma(y_t-y_t^\ast)-\pi_t^e+\varepsilon_t^\pi$
\task $\pi_t=\pi_{t-1}+\gamma(y_t-y_t^\ast)$
\task $\pi_t=\pi^\ast+\gamma(y_t-y_t^\ast)$
\end{tasks}
\ExplainChoices{Correct: inflation equals expected inflation plus a gap term plus a cost-push shock.}{Wrong sign on expected inflation; higher expectations should raise inflation, not lower it.}{Wrong: this is a purely backward-looking accelerationist form; not Whelan’s baseline here.}{Wrong: would force inflation to equal target regardless of the gap/shocks.}



\Question \textbf{Meaning of $\gamma$.} In $\pi_t=\pi_t^e+\gamma(y_t-y_t^\ast)+\varepsilon_t^\pi$, $\gamma$ measures:
\answer{B}
\begin{tasks}(1)
\task how much inflation expectations respond to shocks
\task how much inflation rises when output exceeds natural output by 1 unit (or 1\%)
\task how much the real rate rises when inflation rises
\task the natural rate of unemployment
\end{tasks}
\ExplainChoices{Wrong: that describes expectation formation, not the slope of the PC.}{Correct: $\gamma$ is the slope—how much inflation rises when output exceeds natural output.}{Wrong: that would be a policy-rule parameter, not Phillips slope.}{Wrong: natural unemployment is a level, not a slope parameter.}



\Question \textbf{Inflationary shock.} A positive $\varepsilon_t^\pi$ is best interpreted as:
\answer{C}
\begin{tasks}(1)
\task a permanent increase in productivity
\task a demand boom from fiscal expansion
\task a temporary supply/cost-push shock that raises inflation
\task a fall in expected inflation
\end{tasks}
\ExplainChoices{Wrong: productivity shocks shift natural output, not directly create cost-push inflation.}{Wrong: demand booms are captured by IS shocks; they move output and inflation via the gap.}{Correct: $\varepsilon_t^\pi>0$ is a cost-push/supply shock raising inflation for given activity.}{Wrong: lower expected inflation would reduce inflation, not raise it.}



\Question \textbf{Shift of the Phillips curve.} Holding $y_t-y_t^\ast$ fixed, higher $\pi_t^e$ implies:
\answer{A}
\begin{tasks}(1)
\task higher $\pi_t$ one-for-one (the Phillips curve shifts up)
\task lower $\pi_t$ one-for-one (the Phillips curve shifts down)
\task no change in $\pi_t$
\task ambiguous effect depending on $\gamma$
\end{tasks}
\ExplainChoices{Correct: higher $\pi_t^e$ shifts the Phillips curve up one-for-one in this linear form.}{Wrong: higher expected inflation does not shift PC down.}{Wrong: holding the gap fixed, expectations do change inflation.}{Wrong: the effect is not ambiguous here; it’s additive.}



\Question \textbf{Output gap sign.} If $y_t>y_t^\ast$, then the output gap is:
\answer{D}
\begin{tasks}(1)
\task negative
\task zero
\task undefined
\task positive
\end{tasks}
\ExplainChoices{Wrong: $y_t>y_t^*$ means output is above natural, i.e., positive gap.}{Wrong: gap is zero only if $y_t=y_t^*$.}{Wrong: gap is well-defined as $y_t-y_t^*$.}{Correct: $y_t-y_t^*>0$ implies a positive output gap.}



\Question \textbf{Real vs nominal rate.} The real interest rate relevant for spending decisions is approximately:
\answer{B}
\begin{tasks}(1)
\task $r_t=i_t+\pi_t$
\task $r_t=i_t-\pi_t$
\task $r_t=\pi_t-i_t$
\task $r_t=i_t/\pi_t$
\end{tasks}
\ExplainChoices{Wrong: real rate subtracts inflation; it doesn’t add it.}{Correct: $r_t\approx i_t-\pi_t$ (or expected inflation).}{Wrong: sign is reversed.}{Wrong: division is not the Fisher relationship.}



\Question \textbf{Numerical real rate.} If $i_t=5\%$ and $\pi_t=2\%$, the real rate is approximately:
\answer{C}
\begin{tasks}(1)
\task $7\%$
\task $3\%$? (wrong sign)
\task $3\%$
\task $-3\%$
\end{tasks}
\ExplainChoices{Wrong: adding gives the nominal+inflation, not real.}{Wrong: the option text is ambiguous; the correct computation is $5\%-2\%=3\%$.}{Correct: $r\approx i-\pi=5\%-2\%=3\%$.}{Wrong: negative would require inflation to exceed the nominal rate.}



\Question \textbf{IS curve (Whelan form).} The IS relation is:
\answer{A}
\begin{tasks}(1)
\task $y_t=y_t^\ast-\alpha(i_t-\pi_t-r^\ast)+\varepsilon_t^y$
\task $y_t=y_t^\ast+\alpha(i_t-\pi_t-r^\ast)+\varepsilon_t^y$
\task $y_t=y_t^\ast-\alpha(i_t+\pi_t-r^\ast)+\varepsilon_t^y$
\task $y_t=r^\ast-\alpha(i_t-\pi_t-y_t^\ast)+\varepsilon_t^y$
\end{tasks}
\ExplainChoices{Correct: output falls when the real rate $(i-\pi)$ rises above $r^*$, plus a demand shock.}{Wrong sign: higher real rates would raise output, which contradicts intertemporal substitution.}{Wrong: uses $i+\pi$ instead of the real rate.}{Wrong: mixes variables incorrectly; not the IS form.}



\Question \textbf{Meaning of $\alpha$.} In the IS curve, $\alpha>0$ means:
\answer{B}
\begin{tasks}(1)
\task higher real rates increase output
\task higher real rates reduce output (demand falls with $r$)
\task output is independent of interest rates
\task inflation expectations determine output
\end{tasks}
\ExplainChoices{Wrong: the IS slope is negative in real rates.}{Correct: $\alpha>0$ means higher real rates reduce demand/output.}{Wrong: IS explicitly makes output depend on real rates.}{Wrong: expectations enter via real rates and other blocks, not as the IS slope definition.}



\Question \textbf{Demand shock in IS.} A positive $\varepsilon_t^y$ corresponds to:
\answer{D}
\begin{tasks}(1)
\task a cost-push shock in prices
\task a permanent shift in $y_t^\ast$
\task a higher inflation target $\pi^\ast$
\task a demand shock that raises output for a given real rate
\end{tasks}
\ExplainChoices{Wrong: cost-push shocks belong in the Phillips curve block ($\varepsilon_t^\pi$).}{Wrong: a shift in $y_t^*$ is a supply-side change, not an IS shock.}{Wrong: changing the inflation target is a policy/regime change, not IS shock $\varepsilon_t^y$.}{Correct: $\varepsilon_t^y>0$ raises output for given real rates (demand expansion).}



\Question \textbf{MP rule (baseline).} The simple monetary policy rule is:
\answer{C}
\begin{tasks}(1)
\task $i_t=r^\ast+\pi^\ast-\beta_\pi(\pi_t-\pi^\ast)$
\task $i_t=\pi_t+\beta_\pi(\pi_t-\pi^\ast)$
\task $i_t=r^\ast+\pi^\ast+\beta_\pi(\pi_t-\pi^\ast)$
\task $i_t=r^\ast+\pi_t+\beta_\pi(\pi_t-\pi^\ast)$
\end{tasks}
\ExplainChoices{Wrong sign: would cut rates when inflation rises, destabilizing inflation.}{Wrong: missing the natural real rate and target; also not the standard form.}{Correct: baseline rule raises $i$ when $\pi$ exceeds target, with intercept $r^*+\pi^*$.}{Wrong: mixes in current inflation level incorrectly (double-counting).}



\Question \textbf{Interpret $\beta_\pi$.} A larger $\beta_\pi$ means:
\answer{A}
\begin{tasks}(1)
\task the central bank raises the nominal rate more aggressively when inflation rises
\task the Phillips curve becomes steeper
\task potential output grows faster
\task the IS curve becomes vertical
\end{tasks}
\ExplainChoices{Correct: larger $\beta_\pi$ means stronger nominal-rate response to inflation deviations.}{Wrong: Phillips slope is $\gamma$, not $\beta_\pi$.}{Wrong: potential output growth is a supply-side issue.}{Wrong: IS slope depends on $\alpha$ not $\beta_\pi$.}



\Question \textbf{Deriving IS--MP.} Substituting the baseline MP rule into the IS curve yields an IS--MP curve where output depends on:
\answer{B}
\begin{tasks}(1)
\task only $y_t^\ast$
\task inflation deviations $(\pi_t-\pi^\ast)$
\task only the nominal money supply
\task only expected inflation $\pi_t^e$
\end{tasks}
\ExplainChoices{Wrong: $y^*$ is one component, but IS--MP also depends on inflation via policy response.}{Correct: inflation deviations affect the real rate gap, hence output.}{Wrong: IS--MP is derived without money supply (rate-setting central bank).}{Wrong: expected inflation enters in general, but the baseline derivation focuses on $\pi$ deviations via policy.}



\Question \textbf{Slope of IS--MP.} With baseline MP, the IS--MP slope in $(y,\pi)$ space is proportional to:
\answer{D}
\begin{tasks}(1)
\task $+\alpha(\beta_\pi-1)$ always
\task $-\gamma(\beta_\pi-1)$ always
\task $+\alpha\beta_\pi$ always
\task $-\alpha(\beta_\pi-1)$
\end{tasks}
\ExplainChoices{Wrong: sign should be negative if $\beta_\pi>1$; this expression would be positive.}{Wrong: uses $\gamma$ incorrectly; slope comes from IS and MP, not Phillips.}{Wrong: not the right dependence.}{Correct: IS--MP implies $y-y^*=-\alpha(\beta_\pi-1)(\pi-\pi^*)+\text{shock}$, so slope is $-\alpha(\beta_\pi-1)$.}



\Question \textbf{Taylor principle (sign).} If $\beta_\pi>1$, then the IS--MP curve in $(y,\pi)$ space:
\answer{A}
\begin{tasks}(1)
\task slopes down
\task slopes up
\task is vertical
\task is horizontal
\end{tasks}
\ExplainChoices{Correct: if $\beta_\pi>1$, higher inflation raises the real rate and reduces output → downward IS--MP.}{Wrong: upward slope corresponds to $\beta_\pi<1$.}{Wrong: vertical would mean output independent of inflation, which is not implied.}{Wrong: horizontal would mean inflation independent of output; not implied.}



\Question \textbf{Case $\beta_\pi=1$.} If $\beta_\pi=1$ in the baseline MP rule, then for given shocks the real rate term $(i_t-\pi_t-r^\ast)$ equals:
\answer{C}
\begin{tasks}(1)
\task $(\pi_t-\pi^\ast)$
\task $-(\pi_t-\pi^\ast)$
\task $0$
\task $r^\ast$
\end{tasks}
\ExplainChoices{Wrong: the real-rate gap does not equal inflation deviation in this case.}{Wrong: sign and magnitude are incorrect.}{Correct: if $\beta_\pi=1$, then $i=r^*+\pi$ so $(i-\pi-r^*)=0$.}{Wrong: $r^*$ is the natural rate level, not the gap.}



\Question \textbf{Graph interpretation.} A rightward shift of the IS--MP curve is most naturally associated with:
\answer{B}
\begin{tasks}(1)
\task a negative supply shock ($\varepsilon_t^\pi>0$)
\task a positive demand shock ($\varepsilon_t^y>0$)
\task an increase in expected inflation $\pi_t^e$
\task a decrease in potential output $y_t^\ast$
\end{tasks}
\ExplainChoices{Wrong: that is a supply shock shifting the Phillips curve.}{Correct: $\varepsilon_t^y>0$ raises output at given inflation → IS--MP shifts right.}{Wrong: higher $\pi^e$ shifts PC, not IS--MP (in this model’s decomposition).}{Wrong: lower $y^*$ shifts both blocks but does not correspond to a demand-driven right shift.}



\vspace{0.6em}
\Hrule
\vspace{0.8em}

\section*{Problem 2 (Questions 21--40): Solving IS--MP--PC, expectations, and the Taylor principle}

\Question \textbf{One-period equilibrium.} In the IS--MP--PC model, equilibrium $(y_t,\pi_t)$ is determined by the intersection of:
\answer{D}
\begin{tasks}(1)
\task IS and LM
\task AD and LRAS
\task money demand and money supply
\task IS--MP and PC
\end{tasks}
\ExplainChoices{Wrong: IS-LM is the older model; here we replace LM with MP.}{Wrong: AD-LRAS is not the two-equation closure used here.}{Wrong: money-market clearing is not the closure in IS--MP.}{Correct: equilibrium is where the IS--MP curve intersects the Phillips curve.}



\Question \textbf{Solving strategy.} A standard way to solve the model is to:
\answer{A}
\begin{tasks}(1)
\task substitute IS--MP output gap into the Phillips curve and rearrange
\task substitute the Phillips curve into the MP rule and integrate
\task solve LM for money growth and iterate
\task assume $\pi_t=\pi^\ast$ always
\end{tasks}
\ExplainChoices{Correct: solve by substituting the IS--MP-implied output gap into the Phillips curve.}{Wrong: integrating the MP rule is not the standard solution method here.}{Wrong: LM/money growth is not part of IS--MP--PC.}{Wrong: inflation is not always equal to target with shocks/expectations.}



\Question \textbf{Inflation solution form.} Whelan defines
\[
\theta=\frac{1}{1+\alpha\gamma(\beta_\pi-1)}.
\]
If $\beta_\pi>1$ and $\alpha,\gamma>0$, then $\theta$ satisfies:
\answer{C}
\begin{tasks}(1)
\task $\theta<0$
\task $\theta>1$
\task $0<\theta<1$
\task $\theta=1$
\end{tasks}
\ExplainChoices{Wrong: $\theta$ is positive; not negative.}{Wrong: $\theta$ is less than 1 when $\beta_\pi>1$.}{Correct: with $\alpha,\gamma>0$ and $\beta_\pi>1$, denominator exceeds 1 so $0<\theta<1$.}{Wrong: $\theta=1$ only if $\alpha\gamma(\beta_\pi-1)=0$.}



\Question \textbf{Inflation solution (structure).} The solution for inflation can be written as:
\answer{B}
\begin{tasks}(1)
\task $\pi_t=\pi_t^e+\pi^\ast+\theta(\cdots)$
\task $\pi_t=\theta\pi_t^e+(1-\theta)\pi^\ast+\theta(\text{shocks})$
\task $\pi_t=(1-\theta)\pi_t^e+\theta\pi^\ast+\theta(\text{shocks})$
\task $\pi_t=\pi^\ast$ always
\end{tasks}
\ExplainChoices{Wrong: structure is not additive in that way; weights must sum to 1 on $\pi^e$ and $\pi^*$.}{Correct: inflation is a convex combination of expected inflation and target plus a shock term.}{Wrong: weights reversed relative to the standard derivation.}{Wrong: would eliminate all fluctuation.}



\Question \textbf{Interpretation of $\theta$.} In the inflation solution, a smaller $\theta$ corresponds to:
\answer{A}
\begin{tasks}(1)
\task stronger stabilization of inflation around target (less weight on $\pi_t^e$ and shocks)
\task more sensitivity of inflation to shocks
\task higher natural output
\task a flatter Phillips curve
\end{tasks}
\ExplainChoices{Correct: smaller $\theta$ means less weight on expectations/shocks and more anchoring to target.}{Wrong: more sensitivity corresponds to larger $\theta$.}{Wrong: natural output is not governed by $\theta$.}{Wrong: Phillips slope is $\gamma$, not $\theta$.}



\Question \textbf{Numerical $\theta$.} Let $\alpha=1$, $\gamma=0.5$, $\beta_\pi=3$. Then
$\theta=\frac{1}{1+\alpha\gamma(\beta_\pi-1)}$ equals:
\answer{B}
\begin{tasks}(1)
\task $1/(1+0.5\cdot 3)=1/2.5=0.4$
\task $1/(1+0.5\cdot 2)=1/2=0.5$
\task $1/(1+1\cdot 0.5\cdot 1)=1/1.5\approx 0.667$
\task $1/(1-0.5\cdot 2)=1/0= \infty$
\end{tasks}
\ExplainChoices{Wrong: uses $(\beta_\pi)$ instead of $(\beta_\pi-1)$.}{Correct: $\theta=1/(1+1\cdot0.5\cdot(3-1))=1/(1+1)=0.5$.}{Wrong: uses $(3-2)$ effectively; incorrect.}{Wrong: denominator is positive here; no division by zero.}



\Question \textbf{Inflation target shift.} In $\pi_t=\theta\pi_t^e+(1-\theta)\pi^\ast+\theta(\text{shocks})$, increasing $\pi^\ast$ by 1 point increases $\pi_t$ by:
\answer{D}
\begin{tasks}(1)
\task $\theta$
\task $1$
\task $0$
\task $1-\theta$
\end{tasks}
\ExplainChoices{Wrong: that’s the weight on expectations, not on the target.}{Wrong: inflation does not rise one-for-one unless $\theta=0$.}{Wrong: target matters unless $\theta=1$.}{Correct: coefficient on $\pi^*$ is $(1-\theta)$.}



\Question \textbf{Demand shock effect (sign).} A positive demand shock $\varepsilon_t^y>0$ tends to:
\answer{A}
\begin{tasks}(1)
\task raise both output and inflation (for $\beta_\pi>1$)
\task raise output but reduce inflation
\task reduce output but raise inflation
\task reduce both output and inflation
\end{tasks}
\ExplainChoices{Correct: demand shock shifts IS--MP right → higher output gap → higher inflation via PC (if $\beta_\pi>1$).}{Wrong: inflation typically rises when output rises along the PC.}{Wrong: cost-push, not demand, gives $\pi\uparrow$ with $y\downarrow$.}{Wrong: demand expansion does not reduce both.}



\Question \textbf{Supply shock effect (sign).} A positive inflation shock $\varepsilon_t^\pi>0$ tends to:
\answer{C}
\begin{tasks}(1)
\task raise output and reduce inflation
\task raise both output and inflation
\task raise inflation and (typically) reduce output
\task reduce both output and inflation
\end{tasks}
\ExplainChoices{Wrong: sign is wrong for inflation.}{Wrong: output usually falls when inflation rises from a cost-push shock because policy tightens.}{Correct: $\varepsilon_t^\pi>0$ shifts PC up → inflation up; output typically down.}{Wrong: supply shock does not reduce inflation.}



\Question \textbf{Adaptive expectations.} Under adaptive expectations, Whelan sets:
\answer{B}
\begin{tasks}(1)
\task $\pi_t^e=\pi^\ast$
\task $\pi_t^e=\pi_{t-1}$
\task $\pi_t^e=\mathbb{E}_t\pi_{t+1}$
\task $\pi_t^e=0$ always
\end{tasks}
\ExplainChoices{Wrong: anchored at target is a different assumption.}{Correct: adaptive expectations set $\pi_t^e=\pi_{t-1}$.}{Wrong: forward-looking is NK-style ($\mathbb{E}_t\pi_{t+1}$).}{Wrong: setting to zero is not adaptive.}



\Question \textbf{Inflation persistence.} Under adaptive expectations, inflation persistence arises because:
\answer{D}
\begin{tasks}(1)
\task the IS curve is vertical
\task the MP rule fixes inflation
\task $y_t^\ast$ is random walk
\task today's inflation affects tomorrow's expectations via $\pi_t^e=\pi_{t-1}$
\end{tasks}
\ExplainChoices{Wrong: IS vertical is unrelated.}{Wrong: MP rule does not mechanically fix inflation.}{Wrong: random walk in $y^*$ is not the mechanism in Slides 3.}{Correct: if $\pi_t^e=\pi_{t-1}$, today’s inflation feeds into tomorrow’s expectations.}



\Question \textbf{Taylor principle (core meaning).} ``Taylor principle'' typically means:
\answer{A}
\begin{tasks}(1)
\task $\beta_\pi>1$ (nominal rate rises more than one-for-one with inflation)
\task $\beta_\pi=1$ (real rate constant)
\task $\beta_\pi<1$ (nominal rate reacts less than inflation)
\task $\beta_\pi=0$ (interest rate fixed)
\end{tasks}
\ExplainChoices{Correct: $\beta_\pi>1$ means the nominal rate rises more than one-for-one with inflation, raising the real rate.}{Wrong: $\beta_\pi=1$ keeps the real rate from responding to inflation deviations.}{Wrong: $\beta_\pi<1$ implies real rate falls when inflation rises.}{Wrong: $\beta_\pi=0$ is an interest-rate peg.}



\Question \textbf{IS--MP slope when Taylor principle fails.} If $\beta_\pi<1$, then $-\alpha(\beta_\pi-1)$ is:
\answer{C}
\begin{tasks}(1)
\task negative (downward IS--MP)
\task zero (flat IS--MP)
\task positive (upward IS--MP)
\task undefined
\end{tasks}
\ExplainChoices{Wrong: if $\beta_\pi<1$, then $(\beta_\pi-1)<0$ so $-\alpha(\beta_\pi-1)>0$.}{Wrong: zero slope corresponds to $\beta_\pi=1$.}{Correct: $-\alpha(\beta_\pi-1)$ becomes positive → upward sloping IS--MP.}{Wrong: it is defined; it just changes sign.}



\Question \textbf{Three cases in Slides 3.} Slides 3 discuss different dynamics depending on $\beta_\pi$. One important boundary is:
\answer{D}
\begin{tasks}(1)
\task $\beta_\pi=2$
\task $\beta_\pi=0$
\task $\beta_\pi=\alpha\gamma$
\task $\beta_\pi=1-\frac{1}{\alpha\gamma}$
\end{tasks}
\ExplainChoices{Wrong: not a key boundary in the slides.}{Wrong: $\beta_\pi=0$ is a peg but not the stability boundary highlighted.}{Wrong: $\alpha\gamma$ alone is not the boundary.}{Correct: the boundary $\beta_\pi=1-1/(\alpha\gamma)$ shows up in the adaptive-expectations stability discussion.}



\Question \textbf{When $\theta>1$.} From $\theta=\frac{1}{1+\alpha\gamma(\beta_\pi-1)}$, $\theta>1$ occurs when:
\answer{B}
\begin{tasks}(1)
\task $\beta_\pi>1$
\task $1+\alpha\gamma(\beta_\pi-1)<1$ but still positive, i.e. $\beta_\pi<1$
\task $\alpha\gamma=0$
\task $\beta_\pi=1$
\end{tasks}
\ExplainChoices{Wrong: for $\beta_\pi>1$, denominator is bigger than 1 so $\theta<1$.}{Correct: if $\beta_\pi<1$ but not too small, the denominator can be between 0 and 1, making $\theta>1$.}{Wrong: $\alpha\gamma=0$ would make $\theta=1$, not necessarily $>1$.}{Wrong: $\beta_\pi=1$ implies $\theta=1$.}



\Question \textbf{Explosive adaptive expectations.} In Slides 3, the case $1-\frac{1}{\alpha\gamma}<\beta_\pi<1$ is associated with:
\answer{A}
\begin{tasks}(1)
\task upward and sufficiently steep IS--MP, leading to explosive inflation dynamics under adaptive expectations
\task downward IS--MP and unique stable equilibrium
\task no intersection between IS--MP and PC
\task inflation fixed at target
\end{tasks}
\ExplainChoices{Correct: this region implies an upward and sufficiently steep IS--MP leading to unstable dynamics under adaptive expectations.}{Wrong: downward IS--MP with $\beta_\pi>1$ is the stable case.}{Wrong: intersection exists; the issue is dynamic instability.}{Wrong: inflation is not fixed at target in this region.}



\Question \textbf{Numerical stability check.} Let $\alpha=1$, $\gamma=0.5$. Then $1-\frac{1}{\alpha\gamma}$ equals:
\answer{C}
\begin{tasks}(1)
\task $1-1=0$
\task $1-1/0.5=0.5$
\task $1-2=-1$
\task $1+2=3$
\end{tasks}
\ExplainChoices{Wrong: $1-1/(0.5)$ is not zero.}{Wrong: $1-2$ is $-1$, not $0.5$.}{Correct: $\alpha\gamma=0.5$ so $1-1/0.5=1-2=-1$.}{Wrong: sign is wrong.}



\Question \textbf{Interpretation of $\beta_\pi<1$.} If $\beta_\pi<1$ and inflation rises, then the real rate $i_t-\pi_t$ tends to:
\answer{D}
\begin{tasks}(1)
\task rise strongly, reducing demand
\task stay constant
\task become undefined
\task fall (policy fails to raise the real rate enough)
\end{tasks}
\ExplainChoices{Wrong: policy does not raise the real rate enough.}{Wrong: real rate is not constant when $\beta_\pi<1$.}{Wrong: it is defined.}{Correct: if $i$ rises less than inflation, $i-\pi$ falls → stimulative in real terms.}



\Question \textbf{Shift in expected inflation (graph).} If $\pi_t^e$ increases, the Phillips curve:
\answer{B}
\begin{tasks}(1)
\task shifts down
\task shifts up
\task becomes flatter
\task becomes vertical
\end{tasks}
\ExplainChoices{Wrong: higher expected inflation raises the intercept of the PC.}{Correct: PC shifts up one-for-one in the expectations-augmented form.}{Wrong: slope does not change here; intercept shifts.}{Wrong: PC is not vertical.}



\Question \textbf{Target inflation increase (MP).} In the baseline MP rule $i_t=r^\ast+\pi^\ast+\beta_\pi(\pi_t-\pi^\ast)$, raising $\pi^\ast$ (holding $\pi_t$ fixed) tends to:
\answer{A}
\begin{tasks}(1)
\task lower $i_t$ (because $(\pi_t-\pi^\ast)$ falls)
\task raise $i_t$ one-for-one
\task not change $i_t$
\task raise $i_t$ only if $\beta_\pi>1$
\end{tasks}
\ExplainChoices{Correct: holding $\pi_t$ fixed, raising $\pi^*$ lowers $(\pi_t-\pi^*)$ so the rule sets a lower $i_t$.}{Wrong: not one-for-one; depends on $\beta_\pi$ and the deviation term.}{Wrong: $i$ changes because the deviation term changes.}{Wrong: this effect occurs regardless of $\beta_\pi>1$; it’s algebraic.}



\Question \textbf{Demand shock reversal story.} In Slides 2, a temporary demand shock followed by a reversal can lead to recession with high inflation if:
\answer{C}
\begin{tasks}(1)
\task inflation expectations are perfectly anchored at target
\task the Phillips curve is vertical
\task inflation expectations have drifted up during the boom (adaptive expectations)
\task the Taylor principle is strongly satisfied
\end{tasks}
\ExplainChoices{Wrong: if expectations are anchored, inflation returns faster and the story is weaker.}{Wrong: vertical PC would remove the output-inflation tradeoff; not the point.}{Correct: if expectations drift up during the boom, reversing demand can create recession with still-high inflation.}{Wrong: strong Taylor principle helps stabilize rather than create this pattern.}



\Question \textbf{Soft vs tough central bank.} A ``tough'' central bank in Slides 2 corresponds to:
\answer{B}
\begin{tasks}(1)
\task lower $\beta_\pi$
\task higher $\beta_\pi$
\task higher $\gamma$
\task higher $\varepsilon_t^\pi$
\end{tasks}
\ExplainChoices{Wrong: lower $\beta_\pi$ is ``soft'' (accommodative).}{Correct: higher $\beta_\pi$ means tougher response to inflation.}{Wrong: $\gamma$ is Phillips slope, not central bank toughness.}{Wrong: shock size is not policy toughness.}



\Question \textbf{Effect of higher $\beta_\pi$ on $\theta$.} With $\beta_\pi>1$, increasing $\beta_\pi$ makes $\theta$:
\answer{A}
\begin{tasks}(1)
\task smaller
\task larger
\task unchanged
\task negative
\end{tasks}
\ExplainChoices{Correct: higher $\beta_\pi$ increases the denominator $1+\alpha\gamma(\beta_\pi-1)$, so $\theta$ falls.}{Wrong: it does not rise when $\beta_\pi>1$.}{Wrong: it changes through the denominator.}{Wrong: $\theta$ stays positive in the stable region.}



\Question \textbf{Interpretation: smaller $\theta$.} A smaller $\theta$ implies inflation is:
\answer{D}
\begin{tasks}(1)
\task more sensitive to expectations
\task more sensitive to shocks
\task less tied to the target
\task more tightly anchored to the target
\end{tasks}
\ExplainChoices{Wrong: smaller $\theta$ reduces weight on expectations.}{Wrong: smaller $\theta$ reduces weight on shocks.}{Wrong: smaller $\theta$ makes target more important, not less.}{Correct: smaller $\theta$ means inflation is more tightly anchored to the target.}



\Question \textbf{Compute $\theta$ (quick).} If $\alpha\gamma(\beta_\pi-1)=3$, then $\theta$ equals:
\answer{B}
\begin{tasks}(1)
\task $3$
\task $1/4$
\task $1/3$
\task $4$
\end{tasks}
\ExplainChoices{Wrong: $\theta$ is a fraction, not the denominator.}{Correct: if $\alpha\gamma(\beta_\pi-1)=3$, then $\theta=1/(1+3)=1/4$.}{Wrong: $1/3$ would require denominator 3.}{Wrong: $4$ is the denominator, not $\theta$.}



\Question \textbf{Interpretation of IS--MP intercept.} In the IS equation
$y_t=y_t^\ast-\alpha(i_t-\pi_t-r^\ast)+\varepsilon_t^y$, a higher $r^\ast$ (holding $i_t-\pi_t$ fixed) tends to:
\answer{C}
\begin{tasks}(1)
\task lower output
\task not affect output
\task raise output (real rate gap falls)
\task raise inflation directly
\end{tasks}
\ExplainChoices{Wrong: higher $r^*$ lowers the real-rate gap for given $i-\pi$, so it boosts demand, not lowers it.}{Wrong: it affects output through the real-rate gap term.}{Correct: with fixed $i-\pi$, higher $r^*$ makes $(i-\pi-r^*)$ smaller → higher output.}{Wrong: $r^*$ does not directly shift inflation without going through the gap and PC.}



\Question \textbf{Qualitative: inflation-output loops.} Slides 2 show ``inflation-output loops'' in US data. The point of these figures is mainly to illustrate:
\answer{A}
\begin{tasks}(1)
\task how inflation and output can move in cycles depending on shocks, expectations, and policy
\task that inflation is constant over time
\task that output never deviates from natural output
\task that monetary policy has no effect
\end{tasks}
\ExplainChoices{Correct: the figures illustrate cyclical interactions between activity, inflation, expectations, and policy.}{Wrong: inflation is not constant in the data.}{Wrong: output does deviate from natural output in the short run.}{Wrong: monetary policy affects demand and inflation in the model and (often) in data.}



\end{Exercise}

\end{document}
